\section{Wstęp}

\begin{frame}{Plan prezentacji}
\tableofcontents
\end{frame}

\begin{frame}{Motywacja wyboru tematu}
\begin{itemize}
    \item Uczenie maszynowe poprzez wzmocnienie (Reinforcement Learning, RL) -- dziedzina AI zajmująca się projektowaniem decydentów uczących się przez interakcję z otoczeniem
    \vspace{0.5em}
    \item Tradycyjne podejście: wykładnicze dyskontowanie
    \begin{itemize}
        \item Spójne czasowo
        \item Nie oddaje wielu rzeczywistych zachowań ludzkich
    \end{itemize}
    \vspace{0.5em}
    \item \textbf{Dyskontowanie quasi-hiperboliczne (QH)}
    \begin{itemize}
        \item Bardziej realistyczne modelowanie ludzkich preferencji czasowych
        \item Pozwala na modelowanie niespójności czasowej
        \item Zastosowania: ekonomia behawioralna, psychologia uzależnień, teoria gier
    \end{itemize}
\end{itemize}
\end{frame}

\begin{frame}{Wprowadzenie oznaczeń matematycznych}
\textbf{Markowski Proces Decyzyjny (MDP):} $M = (S, A, q, r, \alpha, \beta)$
\begin{itemize}
    \item $S$ -- zbiór stanów
    \item $A$ -- zbiór akcji
    \item $q(s' \mid s, a)$ -- prawdopodobieństwo przejścia do stanu $s'$ przy akcji $a$ w stanie $s$
    \item $r(s, a)$ -- funkcja nagrody
    \item $\alpha$ -- parametr uprzedzenia teraźniejszości
    \item $\beta$ -- współczynnik dyskontowania wykładniczego
\end{itemize}

\vspace{1em}
\textbf{Funkcja dyskontująca:}
\[
d(t) = 
\begin{cases}
    1, & \text{jeśli } t = 0, \\
    \alpha\,\beta^{t}, & \text{jeśli } t \ge 1.
\end{cases}
\]

\textbf{Skumulowana nagroda:} $R = \sum_{t=0}^{n} d(t)\, r_t$
\end{frame}

\begin{frame}{Wprowadzenie oznaczeń (cd.)}
\textbf{Typy decydentów w zależności od $\alpha$:}
\begin{itemize}
    \item $\alpha = 1$ -- standardowy decydent (czasowo spójny)
    \item $0 < \alpha < 1$ -- decydent quasi-hiperboliczny (czasowo niespójny)
    \item $\alpha = 0$ -- decydent krótkowzroczny
    \item $\alpha > 1$ -- decydent długowzroczny
\end{itemize}

\vspace{1em}
\textbf{Polityka:}
\begin{itemize}
    \item $\phi$ -- polityka (strategia wyboru akcji)
    \item $\mu$ -- reguła pierwszego kroku
    \item $\phi_s$ -- polityka stacjonarna (niezależna od czasu)
    \item $J_{\phi}^{\alpha,\beta}(s)$ -- funkcja wartości dla polityki $\phi$ przy dyskontowaniu QH
\end{itemize}
\end{frame}

\begin{frame}{Kluczowe wyniki -- Twierdzenie o redukcji polityk}
\begin{block}{Twierdzenie 1 \cite{eshwar2025}}
Dla MDP o skończonych zbiorach stanów i akcji, dla dowolnej polityki $\phi = (\phi_0, \phi_1, \ldots)$ istnieje reguła pierwszego kroku $\mu$ oraz stacjonarna polityka $\phi_s$ takie, że:
\[
J_{\phi}^{\alpha,\beta}(s) = J_{(\mu,\phi_s)}^{\alpha,\beta}(s), \qquad \forall s \in S.
\]
\end{block}

\vspace{0.5em}
\textbf{Wniosek:} Wystarczy szukać optymalnej polityki w postaci pary $(\mu^\star, \phi_s^\star)$.

\vspace{0.5em}
\small{Dowód: Jaśkiewicz \& Nowak (2021), Eshwar et al. (2025)}
\end{frame}

\begin{frame}{Kluczowe wyniki -- Operator Bellmana}
\begin{block}{Lemat o kontrakcji \cite{puterman2014,Kallenberg2022}}
Operator Bellmana $T^{\phi_s}$ dla polityki stacjonarnej $\phi_s$ jest kontrakcją z faktorem $\beta$:
\[
\|T^{\phi_s}(V) - T^{\phi_s}(W)\|_\infty \le \beta \|V - W\|_\infty
\]
\end{block}

\vspace{0.5em}
\textbf{Implikacje:}
\begin{itemize}
    \item Gwarancja zbieżności iteracji wartości
    \item Wykładnicza szybkość zbieżności
    \item Możliwość efektywnego obliczania wartości optymalnej polityki
\end{itemize}

\vspace{0.5em}
\small{Odniesienie do dowodu: Puterman (2014), Kallenberg (2022)}
\end{frame}

\begin{frame}{Kluczowe wyniki -- Istnienie optymalnej polityki}
\begin{block}{Twierdzenie 2 \cite{eshwar2025}}
Dla MDP o skończonych zbiorach stanów i akcji istnieje deterministyczna optymalna para $(\mu^\star, \phi_s^\star)$, która składa się z reguły decyzyjnej dla pierwszego kroku i polityki stacjonarnej.
\end{block}

\vspace{0.5em}
\textbf{Znaczenie:}
\begin{itemize}
    \item Redukcja przestrzeni poszukiwań
    \item Możliwość zastosowania algorytmów dyskretnych
    \item Gwarancja znalezienia optymalnego rozwiązania
\end{itemize}

\vspace{0.5em}
\small{Dowód: Eshwar et al. (2025)}
\end{frame}

\begin{frame}{Kluczowe wyniki -- Algorytmy uczenia ze wzmocnieniem}
\textbf{Przedstawione algorytmy:}

\begin{enumerate}
    \item \textbf{QH Policy Evaluation} -- ocena wartości danej polityki
    \begin{itemize}
        \item Aproksymacja stochastyczna o dwóch skalach czasowych
        \item Wykorzystuje warunki Robbinsa-Monro
    \end{itemize}
    
    \vspace{0.5em}
    \item \textbf{QH Q-Learning} -- uczenie optymalnej polityki
    \begin{itemize}
        \item Algorytm bezmodelowy (model-free)
        \item Nie wymaga znajomości dynamiki środowiska
        \item Zbieżność prawie na pewno do optymalnej polityki
    \end{itemize}
\end{enumerate}

\vspace{0.5em}
\begin{block}{Twierdzenie o zbieżności \cite{eshwar2025,borkar2008}}
Algorytm QH Q-Learning zbiega prawie na pewno do optymalnych polityk przy odpowiednich warunkach na kroki uczenia $\eta_n$ i $\theta_n$.
\end{block}

\vspace{0.5em}
\small{Dowód: Eshwar et al. (2025), Borkar (2008)}
\end{frame}
